\section{Validation algorithms} \label{sec:algorithms}

 The previous definitions characterize the correct semantics of a schema
 declaratively as the union of correct typings that a stratification
 happens to pick out, but say nothing about how to compute it.
 The paper
 gives two validation algorithms and states their correctness with respect
 to that declarative semantics:
 Algorithm 1 (\emph{refine}), which computes
 the whole typing \lean{MTyping G S str} at once,
 and Algorithm 2
 (\emph{prove}), which decides membership of a single node-label pair
 without ever constructing the whole typing. Both algorithms, and both
 correctness statements, have been fully formalized and machine-checked in
 Lean.

\subsection{Algorithm 1: refine}

\emph{Refine} computes one stratum at a time: for stratum $i$, it starts
 from the largest plausible typing (everything already fixed for stratum
 $i-1$, plus every possible guess for the new labels of stratum $i$), then
 repeatedly strips out any node-type association that currently fails
 \lean{Satisfies}, until nothing more can be removed --- a greatest-fixed-point
 computation, one per stratum, chained together.

\begin{leanlisting}
def Algorithm1 (G : Graph) (S : Schema) (str : Strat S) : Typing :=
  refineAlg S str.str G str.k

theorem algorithm1_correct (G : Graph) (S : Schema) (str : Strat S)
    (hClosed : S.Closed) : Algorithm1 G S str = MTyping G S str
\end{leanlisting}

The Lean proof mirrors this stratum-by-stratum structure directly. A
single refinement pass (\lean{refineStep}) is iterated to a fixed point
(\lean{refine1}), and a core lemma (\lean{refine1\_eq\_maxAgree}) shows that
this iterative fixed point computes exactly the same typing as
\lean{maxAgree} --- the declarative, union-of-correct-typings construction
that \lean{typ}'s own inductive step uses. Chaining this per-stratum
equivalence across all $k$ strata by induction then gives
\lean{algorithm1\_correct}: two very different-looking constructions,
iterative refinement versus a declarative union, turn out to compute the
same object at every stratum.

\subsection{Algorithm 2: prove}

\emph{Prove} decides whether $(n, L)$ belongs to the typing without ever
constructing it, by recursively visiting only the portion of the
dependency graph reachable from $(n, L)$, tracking a stack of
node-label pairs currently being resolved (to detect and trust positive
cycles) and a set \lean{Dep} of pairs it has confirmed along the way.

\begin{leanlisting}
def prove (S : Schema) (G : Graph) (n : RDFNode) (L : Label) : Bool :=
  proveAux S G (G.nodes ×ˢ S.labels) n L ∅

theorem prove_correct (S : Schema) (str : Strat S) (hClosed : S.Closed)
    (G : Graph) (n : RDFNode) (L : Label) (hLlabels : L ∈ S.labels) :
    prove S G n L = true ↔ (n, L) ∈ MTyping G S str
\end{leanlisting}

This proof is structurally the more delicate of the two. \lean{proveAux} is
a \emph{coinductive-style} algorithm --- when it re-encounters a pair already
on the call stack, it simply trusts it, which is exactly what makes
positive recursive cycles terminate --- so relating it to \lean{MTyping}'s
greatest-fixed-point characterization needs a genuine witness typing, not
just an unfolding argument. The proof first strengthens the semantic
relations \lean{Satisfies}/\lean{Matches} to a \emph{neighbor-local} form,
showing that whether $(n, L)$ holds depends on the ambient typing only
through its restriction to $(n, L)$'s own dependency set --- the fact that
lets \lean{prove} get away with tracking \lean{Dep}/\lean{Hyp} instead of a
whole typing. Soundness and completeness are then proved together, by
strong induction on the stratum, as a single strengthened statement
(\lean{MainProp}) relating \lean{proveAux}'s call-stack bookkeeping to
membership in \lean{typ} at that stratum:

\begin{leanlisting}
def MainProp (S : Schema) (str : Strat S) (G : Graph) (U : Finset Assoc)
    (i : ℕ) : Prop :=
  ∀ Hyp : Finset Assoc,
    (∀ p ∈ Hyp, p.2 ∈ S.stratum str.str i → p ∈ typ S str.str G i) →
    ∀ n L, str.str L = i → L ∈ S.labels →
      (proveAux S G U n L Hyp = true ↔ (n, L) ∈ typ S str.str G i)
\end{leanlisting}

The full argument runs several hundred lines in \lean{Algorithm2.lean};
we refer to it, and to \S4.2 of the paper for the algorithm's original
presentation, rather than reproduce it here.
